\subsection{Motivation}
\label{sec:problem}
\label{sec:setting}
\label{sec:diagnosis}

Consider two episodes in which the agent pushes the same unfamiliar object toward different goals. The first may reveal how an off-center push rotates the object. That observation can help the agent choose a push in the second episode, even when it approaches the object from a different side. An episodic learner corrects its predictions along the first trajectory~\cite{wang2026adajepa}. Carrying its weights forward assumes that the next trajectory calls for a similar correction.

On a stream of PushObj goals, starting each episode from the preceding episode's final weights solved fewer goals than resetting the model (Appendix~D). Consecutive weight updates also point in nearly orthogonal directions (\supref{D}). A trajectory offers a different way to retain what the agent learned. Each executed action and its observed outcome provide a training example that can be revisited alongside contacts from later episodes.

\ours keeps a fast model for the current episode and trains a slow predictor on trajectories from multiple encounters (Fig.~\ref{fig:architecture}). The agent receives no labels telling it when a shape or appearance recurs. It groups trajectories by how they change the model and how well a common predictor explains them. When a new goal arrives, the first few actions provide evidence for which predictor to use.
